\subsection{Scene IV --- The Painting Tool}\label{subsec:scene_paint}

The fourth scene, \texttt{Mano\_Paint\_3D}, is the creative application
(\autoref{subsubsec:sota_creative}). \texttt{LaserPainter} casts a
ray from the index fingertip onto a flat surface --- referred to
throughout as the paper --- and projects the strike onto
the plane of that surface, clamped to its edges. Holding \texttt{Pointing\_Up}
while aimed at the paper lays down ink: the strike point is smoothed to suppress
tracking jitter and lifted a fraction above the surface, and a stroke is grown as a
line whose points are added only once the hand has moved far enough to warrant a new
one. A beam from finger to paper is shown at all times, so the user always sees
where a stroke would begin.

The painting menu, \texttt{BraceletMenuPaint}, offers three controls that are not
all alike. Erase is a mode: it switches the same fingertip from depositing ink to
removing any stroke it passes within a small radius. Save and the lathe toggle are
momentary actions rather than modes, performed once on selection and leaving the
painting state unchanged. The default, with no option active, is to paint.

\begin{figure}[H]
	\centering
	\captionsetup{justification=centering}
	\includegraphics[width=0.9\linewidth, frame]{img/scene_paint_paper.png}
	\caption{Drawing on the paper plane
	         with the fingertip laser.}
	\label{fig:paint-paper}
\end{figure}

Erasing is paired with history. Every drawn and erased stroke is pushed onto an undo
stack, and while erase mode is active the right hand's \texttt{Thumb\_Down} and
\texttt{Thumb\_Up} walk backward and forward through that history, restoring or
removing strokes in turn, with brief on-screen feedback near the wrist. A new stroke
clears the redo branch, as expected of a linear undo history.

The tool is not confined to the plane. \texttt{PainterTraversal}, entered by both
hands' \texttt{Thumb\_Down}, attaches the paper to the player rig and makes it
translucent, then lets the user fly around the strokes already laid down --- both
fists rotating the view in ninety-degree steps, a single fist translating it --- so
that a drawing built up over several viewpoints accumulates into a three-dimensional
form rather than a flat image. A second departure from the plane is the lathe.
Toggling it shrinks the paper to half its width, shifts it aside, and draws an axis
along its new edge; a profile drawn against that axis is treated as the silhouette
of a solid of revolution. The two off-plane modes are mutually exclusive, the lathe
being disabled whenever the three-dimensional mode is entered.

\begin{figure}[H]
	\centering
	\captionsetup{justification=centering}
	\includegraphics[width=0.9\linewidth, frame]{img/scene_paint_lathe.png}
	\caption{Lathe mode: a pawn profile drawn against the
	         rotation axis along the paper's edge.}
	\label{fig:paint-lathe}
\end{figure}

Saving writes the drawing out in a form that depends on how it was made. A plain
two-dimensional drawing is rendered to an image through a temporary orthographic
camera. A three-dimensional drawing is exported as a point cloud, the vertices of
all its strokes written to a PLY file. A lathe profile is first revolved: each
stroke is copied around the axis in many angular steps to sweep out the surface of
the solid it describes, and the resulting points are written as a point cloud in the
same format. The export carries positions only; the per-point normals that surface
reconstruction requires are estimated downstream, exactly as~\autoref{subsec:nksr}
describes for both the Poisson baseline and NKSR.

The exported point cloud then feeds the surface-reconstruction
stage of~\autoref{subsec:nksr}: the three-dimensional and revolved drawings are
exactly the discrete, unstructured samples that PSR and NKSR turn into a mesh,
completing a path from a hand drawn in front of a webcam to a reconstructed
three-dimensional object.

\begin{figure}[H]
	\centering
	\captionsetup{justification=centering}
	\includegraphics[width=0.9\linewidth, frame]{img/scene_paint_reconstruction.png}
	\caption{A mesh reconstructed from a point cloud
	         exported by the painting tool.}
	\label{fig:paint-reconstruction}
\end{figure}
